\section{题目与取舍}

这道题要求为宠物对战游戏做一个 AI 教练，具备军师、陪练、老师三种角色感，参考方向列了 Agentic RL、RAG、ReAct、Memory Stream / Reflection，产品参考是和平精英的「小田」与王者荣耀的「灵宝」。我的理解是：交付物是教练，游戏只是产生局面的环境；三种角色感不是三种语气，而是三种不同的正确性标准——军师要判断准确且可解释，老师只能依据当时可见的信息，陪练要记得住并且不越界。

因此我做的第一个动作是划边界，而不是写代码。三件事本次不做：\textbf{美术与完整游戏内容}，规模控制在 14 只宠物、7 种属性、5 个关卡，这个量已经足够产生「抢先收尾还是留能量」「换宠承伤还是原地治疗」这类真实取舍，继续扩充只会让平衡校准吃掉全部时间；\textbf{联网 PVP}，本地对战是同机分屏轮流操作，服务端仍然接受客户端传来的状态快照，不是权威对局服务，这条差距是架构问题而不是工程量问题；\textbf{模型权重训练}，本机没有可用的 CUDA 后端。

第二个决定与题目「至少完整闭环一项」的要求相反：三种角色我全部做完了。理由是它们的难点并不重叠——军师的难点在分工，老师的难点在时间口径，陪练的难点在克制；只做一种，后两类问题不会暴露出来。
